\begin{abstract}
Los asistentes de domótica basados en modelos de lenguaje suelen depender de
servicios en la nube, lo que introduce alta latencia, costos recurrentes y
dependencia de conectividad. Los Small Language Models (SLMs) ejecutados
localmente son una alternativa atractiva para hubs domóticos de bajo costo,
pero su capacidad real para interpretar comandos de voz en español
rioplatense no ha sido evaluada en profundidad. Este trabajo presenta un
dataset de 32 comandos de domótica en español, con anotaciones de referencia
(intención, dispositivo, ubicación, valor y unidad), y una evaluación
empírica de cuatro SLMs open-source (SmolLM2-360M-Instruct,
Qwen2.5-0.5B-Instruct, Qwen2.5-1.5B-Instruct y SmolLM2-1.7B-Instruct)
ejecutados en CPU con decodificación determinista. Los resultados muestran
una mejora consistente de la coincidencia exacta con el tamaño del modelo
(18,8\% a 59,4\%), a costa de un crecimiento de casi 4x en la latencia de
inferencia (15,6s a 61,9s por comando). La exactitud por campo revela que la
interpretación de valores y unidades numéricas se estanca entre 75\% y 90\%
independientemente del tamaño, y que los cuatro modelos fallan
sistemáticamente (0\%-17\% de exactitud) al interpretar comandos con ajustes
relativos o cualitativos sin un valor explícito, tendiendo a alucinar un
valor o unidad inexistente en el comando original.

\keywords{modelos de lenguaje pequeños, domótica, procesamiento del lenguaje
natural, interpretación de intención, computación en el borde, español
rioplatense.}
\end{abstract}
